\documentclass[11pt]{article}
\usepackage[margin=1in]{geometry}
\usepackage{amsmath,amssymb}
\usepackage{enumitem}

\begin{document}

\section*{Review: \textit{SENA: Sparse-stock Estimation with Nonstationary Adaptation}}

\subsection*{Overall assessment}
This draft is technically stronger than a typical concept paper and now has clearer treatment of irregular intervals, lead-time latent states, and realized-vs-demand-side quantities. The most important remaining issues are implementation-facing consistency problems in state-change reporting and index/boundary conventions, plus a few high-impact ambiguities that can yield materially different implementations.

\subsection*{Most important technical findings}
\begin{enumerate}[leftmargin=*]

\item \textbf{Definite error: reported inventory-change decomposition is inconsistent with the clamped state transition.}

\textbf{Location:} Inventory transition equations in Sections 1/2 ($x_{k,i}=\max\{0, x_{k-1,i}+R_{k,i}-C_{k,i}+A^{\mathrm{adj}}_{k,i}\}$) versus output decomposition in Section 6 ($\Delta x_{k,i}=R_{k,i}-C_{k,i}+A^{\mathrm{adj}}_{k,i}$).

\textbf{Problem:} The decomposition omits the nonnegativity clamp term. If $x_{k-1,i}+R_{k,i}-C_{k,i}+A^{\mathrm{adj}}_{k,i}<0$, the state update clamps to zero, so the stated equality for $\Delta x_{k,i}$ does not hold.

\textbf{Why it matters:} Driver attribution and diagnostics can be wrong exactly in stressed intervals (large negative adjustments / near-stockout states), which are operationally the most important intervals.

\textbf{Suggested fix:} Either (i) define decomposition on a pre-clamp latent state $x^\star_{k,i}=x_{k-1,i}+R_{k,i}-C_{k,i}+A^{\mathrm{adj}}_{k,i}$, or (ii) include an explicit clamp-correction term in $\Delta x_{k,i}$.

\item \textbf{Likely issue: interval indexing and initial conditions are under-specified.}

\textbf{Location:} Definitions use report times $t_1<\cdots<t_K$ and $\Delta t_k=t_k-t_{k-1}$, while transitions use $x_{k-1,i}$, $B_{k-1,i}$, and indicators on $[t_{k-1},t_k]$.

\textbf{Problem:} The manuscript does not define $t_0$, $x_{0,i}$, or $B_{0,i}$ (or an explicit interval index set such as $k=2,\ldots,K$ for updates).

\textbf{Why it matters:} This creates ambiguity in the first interval and can change likelihood terms, receipt assignment, and initial posterior mass.

\textbf{Suggested fix:} Add an explicit indexing convention (for example, updates for $k=2,\dots,K$ with initial states at $k=1$) or define $t_0$ and priors/initialization for $x_{0,i},B_{0,i}$.

\item \textbf{Likely issue: order-policy coefficients are described as time-varying but are not given time dynamics.}

\textbf{Location:} State list and placement model ($\alpha_{k,i},\beta_{k,i}$ appear time-indexed), plus hyperparameter priors.

\textbf{Problem:} $\alpha_{k,i}$ and $\beta_{k,i}$ are indexed by $k$ but no transition model is provided; priors are given directly, which implies either independent per-interval coefficients or an incomplete specification.

\textbf{Why it matters:} This changes parameter count and identifiability substantially and can make implementations diverge in both complexity and behavior.

\textbf{Suggested fix:} Choose one explicit design: static coefficients ($\alpha_i,\beta_i$) or dynamic coefficients with stated transitions (e.g., random walk with shrinkage).

\item \textbf{Likely implementation ambiguity: LogNormal parameterization is inconsistent across equations.}

\textbf{Location:} Order-size model uses $\mathrm{LogNormal}(\mu^O_{k,i},\sigma^O_i)$, while lead-time model uses $\mathrm{LogNormal}(\mu_{L,u,i},\sigma^2_{L,u,i})$.

\textbf{Problem:} One place appears to use a log-standard-deviation parameter and the other a log-variance parameter.

\textbf{Why it matters:} Different software APIs interpret LogNormal arguments differently; this ambiguity can produce materially different variances and reorder-risk outputs.

\textbf{Suggested fix:} State one global convention explicitly (e.g., ``LogNormal(log-mean, log-standard-deviation)'') and rewrite both formulas to match that convention.

\item \textbf{Unsupported claim: causal language about promo modeling is stronger than evidence shown.}

\textbf{Location:} Promotions motivation text (``prevents'' misattribution and ``improves identifiability'').

\textbf{Problem:} The manuscript proposes an evaluation protocol but does not yet present empirical evidence demonstrating prevention-level causal claims.

\textbf{Why it matters:} Overstated language weakens scientific calibration and may be challenged in review.

\textbf{Suggested fix:} Rephrase to ``can reduce misattribution risk'' and ``can improve identifiability,'' then tie the claim to the evaluation section as a falsifiable hypothesis.

\end{enumerate}

\subsection*{Notation / convention drift check (required)}
I explicitly checked recurring high-risk symbols for descriptor drift (first definition versus later usage), including $x_{k,i}$, $D_{k,i}$, $C_{k,i}$, $B_{k,i}$, $S_{k,j}$, $Q_{k,i}$, $m^L_{k,i}$, and $v^L_{k,i}$.

\begin{itemize}[leftmargin=*]
\item No initial/final, source/observer, or role-adjective conflicts were found for these symbols.
\item Sign conventions for consumption and pipeline terms are internally consistent after the realized-demand clarification.
\item Branch/quadrant checks were performed: no inverse-trig expressions were found, so no concrete atan/atan2 branch issue is present.
\item Dimensional checks on key planning quantities ($d_{k,i}$, $D^{\mathrm{LT}}_{k,i}$, DoC, SS, ROP) did not reveal a concrete unit mismatch.
\end{itemize}

\subsection*{Meaningful editorial findings}
\begin{enumerate}[leftmargin=*]
\item \textbf{Location:} Scope wording in Abstract/Section 1 (``single-business'' versus ``single-user'').

\textbf{Problem:} The manuscript alternates between two scope labels that can be interpreted differently.

\textbf{Why it matters:} Scope ambiguity affects reproducibility and external validity claims.

\textbf{Suggested fix:} Pick one term and define it once (e.g., ``single business entity with one internally coherent inventory process'').

\textbf{Candidate rewrite:} ``SENA is designed for a single business entity with intermittently observed inventory and no cross-business pooling.''

\item \textbf{Location:} Equation-dense sections (inventory transition, lead-time mapping, ROP/SS definitions).

\textbf{Problem:} Many key equations are referenced verbally without equation labels.

\textbf{Why it matters:} This slows technical review and increases risk of implementation mistakes when citing formulas.

\textbf{Suggested fix:} Add equation labels for core transitions and refer to them consistently in prose.
\end{enumerate}

\subsection*{Proofing-sweep (required micro-scan)}
I ran the bundled proofing scan logic on source and compiled outputs.

\begin{itemize}[leftmargin=*]
\item \texttt{main.tex}: no high-confidence scan hits.
\item \texttt{refs.bib}: no high-confidence scan hits.
\item \texttt{main.pdf}: repeated ``double period'' hits arise from table-of-contents dot leaders; spot-check indicates these are formatting artifacts, not punctuation defects in manuscript prose.
\end{itemize}

\subsection*{Highest-priority fixes before submission}
\begin{enumerate}[leftmargin=*]
\item Repair Section 6 inventory-change decomposition so it is mathematically consistent with the nonnegativity clamp.
\item Make interval indexing and initial-state conventions explicit ($t_0/x_0/B_0$ or equivalent $k$ range).
\item Resolve time-indexing intent for order-policy coefficients ($\alpha,\beta$ static vs dynamic with transition equations).
\item Standardize LogNormal parameterization notation across order-size and lead-time equations.
\end{enumerate}

\subsection*{Items needing external verification}
No high-priority findings above require outside literature lookup; they are internal consistency and implementation-specification issues within the manuscript.

\end{document}